\documentclass{article}

% if you need to pass options to natbib, use, e.g.:
%     \PassOptionsToPackage{numbers, compress}{natbib}
% before loading neurips_2026

\usepackage{neurips_2026}

\usepackage[utf8]{inputenc} % allow utf-8 input
\usepackage[T1]{fontenc}    % use 8-bit T1 fonts
\usepackage{hyperref}       % hyperlinks
\usepackage{url}            % simple URL typesetting
\usepackage{booktabs}       % professional-quality tables
\usepackage{amsfonts}       % blackboard math symbols
\usepackage{nicefrac}       % compact symbols for 1/2, etc.
\usepackage{microtype}      % microtypography
\usepackage{xcolor}         % colors
\usepackage{multirow}       % \multirow in tables
\usepackage{amsmath}        % math formatting
\usepackage{natbib}
\usepackage[section]{placeins}  % keeps floats within their section

\title{When Does One-Shot Pruning Beat Iterative Optimisation?\\
       Second-Order Correction for Sparse LLMs on Neuromorphic Hardware}

\author{%
  David S.~Hippocampus\thanks{Use footnote for providing further information
    about author (webpage, alternative address)---\emph{not} for acknowledging
    funding agencies.} \\
  Department of Computer Science\\
  Cranberry-Lemon University\\
  Pittsburgh, PA 15213 \\
  \texttt{hippo@cs.cranberry-lemon.edu} \\
}

\begin{document}

\maketitle

\begin{abstract}
Post-training weight pruning for large language models must solve two coupled problems: selecting which weights to remove and correcting the remaining weights to compensate. We study when one-shot second-order correction beats iterative activation-aware optimisation, and vice versa. We propose \textbf{OBS-cancel-block}, a greedy pruning criterion that identifies \emph{cancellation groups}---sets of correlated weights whose joint output contribution is smaller than any per-weight score predicts---and corrects the remaining weights via Optimal Brain Surgeon residuals within fixed-width column blocks. Evaluating across three model families (1B FLA transformer, LLaMA-7B, HGRN-1.3B SSM) and sparsity levels from 30\% to 80\%, we uncover a clear phase transition: at 1B--1.3B scale, OBS-cancel-block dominates all baselines including AWP, an iterative hard-thresholding method optimising the same objective; at 7B scale, AWP's 200-iteration search finds better sparse solutions at a cost OBS-cancel-block does not require. For neuromorphic deployment on Intel Loihi---where model scale is constrained by hardware memory---OBS-cancel-block delivers the best accuracy-efficiency tradeoff: it outperforms SparseGPT by up to $1.83\times$ at 80\% sparsity, generalises across both attention and gated recurrent architectures, and runs in a single forward pass. To our knowledge, this is the first systematic study of billion-parameter sparse LLM inference on neuromorphic hardware.
\end{abstract}


% =====================================================================
\section{Introduction}
\label{sec:intro}
% =====================================================================

Post-training weight pruning compresses large language models without retraining by removing weights and optionally correcting the survivors. Two broad paradigms have emerged: \emph{one-shot} methods that score and select in a single pass, and \emph{iterative} methods that alternate between sparse projection and gradient steps. Both optimise variants of the same layer-wise reconstruction objective, but they differ fundamentally in how much search they perform---and in how much compute they consume.

The question of when each paradigm wins has not been studied systematically. We close this gap. Our setting is weight pruning for inference on Intel Loihi, a neuromorphic chip whose event-driven, memory-constrained architecture makes weight sparsity uniquely valuable: zero weights generate no spike events, no memory reads, and no switching activity, so sparsity translates to proportional savings without the structured-format constraints that limit unstructured pruning on GPUs. Loihi operates comfortably in the 1B--2B parameter regime; this makes understanding pruning at \emph{that} scale the practically relevant question.

We propose \textbf{OBS-cancel-block}, a one-shot second-order method that extends Optimal Brain Surgeon (OBS) to identify \emph{cancellation groups}: sets of correlated input channels whose weights carry opposite signs, so their joint contribution to the layer output is smaller than any per-weight score predicts. We evaluate it alongside Wanda, RIA, SparseGPT, and AWP---an iterative hard-thresholding baseline optimising the same objective as OBS-cancel-block but via up to 200 gradient-projection iterations per layer---across three model families and six sparsity levels (30--80\%).

Our main contributions are:

\begin{itemize}
  \item \textbf{A scale-dependent phase transition in pruning effectiveness.}
        At 1B--1.3B scale, one-shot second-order correction (OBS-cancel-block)
        dominates all methods including AWP: at 50\% sparsity on HGRN-1.3B,
        OBS-cancel-block achieves PPL~19.0 while AWP collapses to 527,
        comparable to uncorrected Wanda (584). At 7B scale the relationship
        inverts: AWP achieves PPL~230 vs.\ 948 for OBS-cancel-block at 80\%
        sparsity on LLaMA-7B. We provide a mechanistic explanation: at
        sub-2B scale, reconstruction error magnitude is the binding constraint
        and second-order correction is the only lever that controls it;
        at 7B scale, the optimisation landscape has sufficient redundancy
        that iterative search finds markedly better sparse solutions.

  \item \textbf{OBS-cancel-block: one-shot, second-order, architecture-agnostic.}
        Our method outperforms SparseGPT at every sparsity level from 30\% to
        80\% on all evaluated models ($1.02\times$ to $1.83\times$) using the
        same weight corrections---the gain comes entirely from the
        cancellation-aware greedy selection criterion. It generalises without
        modification across attention-based transformers and gated recurrent
        SSMs (HGRN), and is preferred over AWP for neuromorphic deployment on
        three grounds: better quality at the relevant scale, a single
        forward-pass cost (${\sim}200\times$ less compute than AWP), and
        no sensitivity to iteration count hyperparameters.

  \item \textbf{First systematic sparse LLM deployment on Intel Loihi.}
        We map pruned 1B and 1.3B models onto Loihi and measure energy per
        token and latency for both prefill and decode workloads. Weight
        sparsity at 50\%+ yields concrete efficiency gains during decode;
        we quantify how algorithmic pruning quality translates to
        hardware efficiency and establish a co-design reference for
        future neuromorphic LLM work.
\end{itemize}


% =====================================================================
\section{Background}
\label{sec:background}
% =====================================================================

\paragraph{Magnitude pruning.}
The earliest systematic pruning methods remove weights with the smallest
absolute value \citep{han2015learning}. This criterion is independent of the
input activation distribution: a small weight connected to a high-variance
input channel can matter far more than a large weight connected to a
near-zero one.

\paragraph{Activation-weighted scoring.}
Wanda \citep{sun2023simple} scores each weight as the product of its magnitude
and the root-mean-square (RMS) activation of the corresponding input channel,
coupling importance to activation scale without requiring a matrix inverse.
RIA \citep{zhang2024plug} extends this by normalising each weight by both its
row and column $\ell_1$ norms, making importance relative to peer weights in
the same output neuron and input channel. Both methods apply no weight
correction after pruning.

\paragraph{Second-order methods.}
SparseGPT \citep{frantar2023sparsegpt} applies Optimal Brain Surgeon (OBS)
saliency \citep{hassibi1993optimal,lecun1989optimal} to score weights by the
second-order reconstruction error they induce, and applies sequential OBS
corrections within blocks of 128 columns to partially compensate for pruning
error. This reduces perplexity degradation substantially at small scales and
high sparsity, at the cost of requiring the full layer Hessian.

\paragraph{Iterative activation-aware methods.}
AWP \citep{liu2025awp} reframes pruning as minimising the same activation-weighted
reconstruction error as OBS methods,
$\|\mathbf{W}\boldsymbol{\Sigma}_X^{1/2} - \hat{\mathbf{W}}\boldsymbol{\Sigma}_X^{1/2}\|_F^2$,
but solves it via Iterative Hard Thresholding (IHT) rather than one-shot greedy
selection. Starting from a Wanda initialisation, each iteration applies a gradient
step with Lipschitz step size $\eta = 2/\|\boldsymbol{\Sigma}_X\|_F$ followed by
projection onto the set of masks with exactly $k$ zeros per row. Convergence
requires up to 200 passes over the layer weight matrix per layer, making AWP
substantially more expensive than single-pass methods but potentially better able to
escape the greedy solution.

\paragraph{What prior methods miss: cancellation groups.}
Wanda, RIA, and SparseGPT score each weight independently or via diagonal
covariance proxies. When two input channels are positively correlated and
their weights carry opposite signs, their joint contribution to the layer
output is smaller than any per-weight score predicts. These \emph{cancellation
groups} are invisible to all diagonal-scoring criteria. AWP uses the full
covariance matrix and so is not limited by diagonal approximation; however,
its IHT-based solver still does not explicitly identify cancellation structure---it
discovers it only insofar as the gradient descent trajectory passes through
configurations where cancellation groups are jointly pruned.

\paragraph{Hardware design space.}
All existing methods were designed for GPU execution models, where throughput
gains from unstructured sparsity require highly structured formats (e.g., 2:4
block sparsity) to avoid irregular memory access. On event-driven neuromorphic
hardware, unstructured zeros are natively efficient---zero-weight synapses
generate no spike events and consume no memory bandwidth regardless of their
position in the weight matrix. This removes the structured-format constraint
and opens the pruning design space to criteria that prioritise reconstruction
quality over pattern regularity.


% =====================================================================
\section{Methods}
\label{sec:methods}
% =====================================================================

\subsection{Problem Formulation}

For each linear layer with weight matrix
$\mathbf{W} \in \mathbb{R}^{d_\text{out} \times d_\text{in}}$, post-training
pruning selects a binary mask
$\mathbf{M} \in \{0,1\}^{d_\text{out} \times d_\text{in}}$ and applies it as
$\hat{\mathbf{W}} = \mathbf{M} \odot \mathbf{W}$.
We adopt semi-structured sparsity: exactly
$k = \lfloor d_\text{in} \cdot s \rfloor$ weights are zeroed per output row,
where $s$ is the target sparsity ratio. This decouples the per-row selection
problem across output neurons while preserving a regular memory layout
compatible with spatial accelerators.

Each output row is treated independently. For row $i$, the pruning objective
is to choose index set $S \subset [d_\text{in}]$ with $|S| = k$ that minimises
the increase in squared output error under the data distribution:
\begin{equation}
  \min_{\substack{S \subset [d_\text{in}] \\ |S| = k}} \;
  \mathbb{E}_x \!\left[
    \Bigl(\mathbf{w}_i^\top x - (\mathbf{M}_i \odot \mathbf{w}_i)^\top x\Bigr)^{\!2}
  \right]
  = \mathbf{w}_{i,S}^\top \,\boldsymbol{\Sigma}_X[S,S]\, \mathbf{w}_{i,S},
  \label{eq:obj}
\end{equation}
where $\boldsymbol{\Sigma}_X = \mathbb{E}[xx^\top]$ is the input activation
covariance and $\mathbf{w}_{i,S}$ are the pruned weights.


\subsection{Activation Statistics}

All methods require statistics of the input activations. We register forward
hooks on every linear layer and collect a calibration set of 64 batches
(batch size 4, sequence length 512) drawn from C4. For methods that only
need the diagonal of $\boldsymbol{\Sigma}_X$ (Wanda, RIA), we accumulate the
per-channel RMS activation:
\begin{equation}
  \hat{\sigma}_j = \sqrt{\frac{1}{N}\sum_{n=1}^{N} x_{n,j}^2}.
  \label{eq:rms}
\end{equation}
For methods requiring the full covariance (OBS-cancel-block, SparseGPT), we
accumulate $\mathbf{H} = \frac{2}{N} \mathbf{X}^\top \mathbf{X}$ (the
empirical Hessian of the squared output error), then compute $\mathbf{H}^{-1}$
via Cholesky decomposition with a damping term
$\lambda = 0.01 \cdot \operatorname{mean}(\operatorname{diag}(\mathbf{H}))$
for numerical stability.


\subsection{Baseline Scoring Methods}

\paragraph{Wanda.}
Scores each weight independently as the product of its absolute value and the
RMS activation of its input channel:
\begin{equation}
  s^\mathrm{Wanda}_{i,j} = |W_{i,j}| \cdot \hat{\sigma}_j.
\end{equation}
Per output row, the $k$ lowest-scoring weights are zeroed. No weight
correction is applied.

\paragraph{RIA.}
Extends Wanda by normalising by both the row and column $\ell_1$ norms:
\begin{equation}
  s^\mathrm{RIA}_{i,j}
  = \left(
      \frac{|W_{i,j}|}{\|\mathbf{w}_i\|_1}
      + \frac{|W_{i,j}|}{\|\mathbf{W}_{:,j}\|_1}
    \right)
    \cdot \hat{\sigma}_j^{\,\alpha},
  \qquad \alpha = 0.5.
\end{equation}
No weight correction is applied.

\paragraph{SparseGPT.}
Uses the OBS saliency score $W_{i,j}^2 / [\mathbf{H}^{-1}]_{jj}$ to select
weights in column order within blocks of 128. After each selected weight is
zeroed, the remaining weights in the block are updated via a rank-1 OBS
correction:
\begin{equation}
  \delta \mathbf{w} = -\frac{W_{i,j}}{[\mathbf{H}^{-1}]_{jj}}
                       \cdot [\mathbf{H}^{-1}]_{:,j},
\end{equation}
and $\mathbf{H}^{-1}$ is updated via a Schur complement step. This correction
redistributes the pruned weight's contribution onto the remaining weights in
the block.


\subsection{Cancellation-Aware Pruning: OBS-cancel-block}
\label{sec:method_ours}

The objective in Equation~\eqref{eq:obj} is a quadratic form in the pruned
weights. When two input channels $j, j'$ are positively correlated
($\boldsymbol{\Sigma}_X[j,j'] > 0$) and their weights carry opposite signs
($W_{i,j} \cdot W_{i,j'} < 0$), their joint contribution to the layer output
is smaller than either per-weight score predicts. Standard diagonal scoring
(Wanda, RIA, SparseGPT) is blind to these cancellation groups.

\paragraph{Greedy selection via OBS residuals.}
We solve the combinatorial selection problem greedily. After pruned set $S'$
has been chosen, the marginal cost of additionally pruning column $j$ is:
\begin{equation}
  \delta(j \mid S') = \frac{r_j^2}{d_j},
\end{equation}
where $r_j$ and $d_j$ are the $j$-th residual and diagonal element of the
evolving Cholesky factor of $\mathbf{H}^{-1}$, updated via rank-1 Schur
complement steps after each selection. At the first step ($S' = \emptyset$)
this recovers the SparseGPT score exactly; at subsequent steps it accounts
for correlations between already-pruned and candidate weights.

\paragraph{Block-level variant (OBS-cancel-block, proposed).}
On models with $d_\text{in} \geq 2048$, global greedy selection breaks down
for two reasons: (1) the selected mask is not column-ordered, violating the
assumption of the subsequent OBS correction; and (2) thousands of Schur
complement rank-1 updates cause the residual diagonal to drift numerically
even in float64. We resolve both by restricting selection and correction to
fixed-width blocks of $B = 128$ columns. Within each block,
$\lfloor B \cdot s \rfloor$ OBS-cancel steps select weights to prune from the
block's $B \times B$ submatrix of $\mathbf{H}^{-1}$; OBS corrections are
applied immediately after. Blocks are processed left to right, eliminating
the ordering mismatch and capping numerical drift at $B$ updates per block.

On the 1B model, where global OBS-cancel is numerically stable, the block
variant captures slightly less cross-block cancellation (PPL 25.5 vs.\ 24.1
at 50\% sparsity); on LLaMA-7B and HGRN-1.3B, it is strictly superior to the
global variant.

\subsection{Evaluation Protocol}

\paragraph{Perplexity.}
Measured on the WikiText-2 test set \citep{merity2016pointer}, token-by-token
with a sliding window of 512 tokens.

\paragraph{Downstream accuracy.}
Zero-shot evaluation using \texttt{lm-evaluation-harness}
\citep{eval-harness} on ARC-Easy, ARC-Challenge, HellaSwag, PIQA,
WinoGrande, and LAMBADA. We report the unweighted average (\textbf{Avg}).

All pruning is one-shot with no fine-tuning; models are loaded in bfloat16.


% =====================================================================
\section{Hardware Deployment on Intel Loihi}
\label{sec:hardware}
% =====================================================================

% -----------------------------------------------------------------------
% TODO (Sumit): Please insert your methods and results in this section.
%
% Suggested subsections and content:
%
%   \subsection{Loihi Architecture and Sparse Execution Model}
%     - Describe the Loihi hardware model: cores, synaptic memory layout,
%       event-driven execution, and how zero weights map to zero spike events.
%     - Explain why semi-structured (constant-k-per-row) sparsity is
%       preferable to fully unstructured sparsity on Loihi (bounded load
%       imbalance across cores).
%
%   \subsection{Mapping Sparse LLM Layers to Loihi}
%     - Describe how pruned linear layers are compiled and mapped to Loihi
%       cores, and how the sparsity format is represented in synaptic memory.
%     - Note any model-specific considerations for the 1B transformer vs.
%       HGRN-1.3B.
%
%   \subsection{Inference Benchmarks: Prefill and Decode}
%     - Energy per token and latency measurements for 50% and 80% sparse
%       models vs. the dense baseline, for both prefill and decode workloads.
%     - Comparison of weight-sparsity gains (decode) vs. activation-sparsity
%       gains (prefill) on Loihi, motivating the unified sparsity strategy.
%     - If available: comparison of OBS-cancel-block vs. SparseGPT sparse
%       models on Loihi hardware (not just PPL), to connect algorithmic
%       quality to hardware efficiency.
%
% Key claims to establish with data:
%   - This is among the first deployments of billion-parameter sparse LLMs
%     on neuromorphic hardware.
%   - Weight sparsity yields proportional reductions in active synaptic
%     operations during decode; activation sparsity is the primary lever
%     during prefill.
%   - Semi-structured sparsity at 50%+ enables meaningful energy and latency
%     reductions without catastrophic accuracy loss (pointing to the
%     PPL/accuracy tables in Section~\ref{sec:experiments}).
% -----------------------------------------------------------------------

\subsection{Loihi Architecture and Sparse Execution Model}

[Sumit: please describe the Loihi hardware model and sparse execution model here.]

\subsection{Mapping Sparse LLM Layers to Loihi}

[Sumit: please describe the compilation and mapping of pruned layers to Loihi cores here.]

\subsection{Inference Benchmarks: Prefill and Decode}

[Sumit: please insert hardware benchmark results here, including energy per token, latency, and comparisons with the dense baseline across prefill and decode workloads.]


% =====================================================================
\section{Experiments}
\label{sec:experiments}
% =====================================================================

\subsection{Experimental Setup}

\paragraph{Models.}
We evaluate on three model families spanning two architectures and two scales:
(1)~a 1B-parameter FLA transformer, (2)~LLaMA-7B
\citep{touvron2023llama} (\texttt{openlm-research/open\_llama\_7b}), and
(3)~HGRN-1.3B \citep{qin2023hierarchically}
(\texttt{fla-hub/hgrn-1.3B-100B}), a gated recurrent SSM with no attention
layers. Including both attention-based and recurrent architectures tests
whether gains from OBS-cancel-block are architecture-agnostic.

\paragraph{Calibration.}
All methods use 64 calibration batches (batch size~4, sequence length~512)
drawn from C4~\citep{raffel2020exploring} (2{,}000 streaming documents).
Per-channel activation RMS values are accumulated via forward hooks on every
linear layer; the full Hessian $\mathbf{H} = \frac{2}{N}\mathbf{X}^\top
\mathbf{X}$ is used by OBS-cancel-block and SparseGPT.

\paragraph{Baselines.}
We compare against Wanda~\citep{sun2023simple}, RIA~\citep{zhang2024plug},
SparseGPT~\citep{frantar2023sparsegpt}, and AWP~\citep{liu2025awp}. We also
include two intermediate ablations: greedy covariance pruning without weight
correction (\textsc{Greedy-cov}), and greedy covariance pruning with closed-form weight
correction (\textsc{Greedy+corr}).

% -----------------------------------------------------------------------
\subsection{Main Results: 1B Transformer at 50\% Sparsity}
\label{sec:main_results}

\begin{table}[htbp]
  \centering
  \caption{%
    1B transformer at 50\% sparsity.
    PPL on WikiText-2; downstream accuracy zero-shot.
    Best non-dense result in \textbf{bold}.
  }
  \label{tab:main_1b}
  \small
  \begin{tabular}{lrrrrrrr}
    \toprule
    Method & PPL$\downarrow$ & ARC-e & ARC-c & Hella. & PIQA & Wino. & Avg$\uparrow$ \\
    \midrule
    Dense baseline       & 19.5  & .610 & .268 & .366 & .683 & .525 & .472 \\
    \midrule
    Wanda                & 3{,}545 & .271 & .228 & .260 & .523 & .493 & .296 \\
    RIA                  & 2{,}105 & .277 & .263 & .268 & .513 & .490 & .302 \\
    SparseGPT            &    29.6 & .555 & .294 & .428 & .665 & .517 & .492 \\
    \midrule
    \textit{Greedy-cov (ablation)}                      & 1{,}103 & .279 & .216 & .267 & .535 & .476 & .355 \\
    \textit{Greedy+corr (ablation)}                     &     615 & .291 & .205 & .272 & .545 & .471 & .298 \\
    \textit{Hybrid: cancellation+OBS (ablation)}        &     958 & .279 & .208 & .265 & .546 & .470 & .354 \\
    \textit{Interleaved: block cancel.+OBS (ablation)}  &     856 & .285 & .238 & .257 & .534 & .492 & .361 \\
    \textit{Global OBS-cancel (ablation, $b\to\infty$)} &  24.1   & .569 & .248 & .345 & .665 & .517 & .469 \\
    \midrule
    \textbf{OBS-cancel-block (ours)}  & \textbf{25.5} & \textbf{.561} & .247 & .345 & \textbf{.662} & .512 & \textbf{.465} \\
    \bottomrule
  \end{tabular}
\end{table}

Table~\ref{tab:main_1b} shows all methods on the 1B transformer at 50\%
sparsity. OBS-cancel-block achieves PPL~25.5, outperforming SparseGPT (29.6)
by \textbf{1.16$\times$} with identical OBS weight corrections; the gain
comes entirely from the cancellation-aware greedy selection criterion.

The ablation rows expose why naive combinations fail. Both the hybrid
(cancellation scores with OBS correction) and interleaved (block-level
cancellation then OBS) variants collapse to PPL\,$>$\,800, far worse than
either pure method. OBS correction is internally consistent with SparseGPT's
score $w^2/[\mathbf{H}^{-1}]_{jj}$; substituting an externally-derived score
breaks this coupling. The fix is to derive cancellation-awareness within the
OBS objective via the greedy $r_j^2/d_j$ criterion, which is exactly what
OBS-cancel-block does.

The global OBS-cancel ablation ($b\to\infty$, PPL~24.1) represents the
theoretical ceiling of the method when block constraints are lifted; it is
only stable on the 1B model (see Section~\ref{sec:llama}).

No-correction methods (Wanda, RIA, Greedy-cov) all collapse to
PPL\,$>\,$600 at 50\% sparsity, confirming that weight correction is
necessary at this scale regardless of importance score quality.

% -----------------------------------------------------------------------
\subsection{Calibration Data Ablation}
\label{sec:calib_ablation}

\begin{table}[htbp]
  \centering
  \caption{%
    Calibration data ablation on the 1B transformer at 50\% sparsity.
    PPL always measured on WikiText-2 test set.
  }
  \label{tab:calib}
  \small
  \begin{tabular}{llrrrrrrr}
    \toprule
    Calib.\ data & Method & PPL$\downarrow$ & ARC-e & ARC-c & Hella. & PIQA & Wino. & Avg$\uparrow$ \\
    \midrule
    \multirow{4}{*}{WikiText-2}
      & Wanda            & 3{,}545 & .271 & .228 & .260 & .523 & .493 & .296 \\
      & RIA              & 2{,}105 & .277 & .263 & .268 & .513 & .490 & .302 \\
      & SparseGPT        &    29.6 & .555 & .294 & .428 & .665 & .517 & .492 \\
      & OBS-cancel-block &    25.5 & .561 & .247 & .345 & .662 & .512 & .465 \\
    \midrule
    \multirow{4}{*}{C4}
      & Wanda            & 4{,}051 & .271 & .275 & .267 & .502 & .489 & .300 \\
      & RIA              & 2{,}357 & .270 & .267 & .272 & .515 & .496 & .303 \\
      & SparseGPT        &    31.0 & .505 & .298 & .429 & .678 & .517 & .435 \\
      & \textbf{OBS-cancel-block} & \textbf{26.3} & \textbf{.519} & .291 & \textbf{.426} & \textbf{.663} & \textbf{.541} & \textbf{.443} \\
    \bottomrule
  \end{tabular}
\end{table}

Under WikiText-2 calibration, SparseGPT has a task-accuracy advantage over
OBS-cancel-block (0.492 vs.\ 0.465). With C4 calibration this reverses:
OBS-cancel-block outperforms SparseGPT on both PPL (26.3 vs.\ 31.0) and
downstream accuracy (0.443 vs.\ 0.435). SparseGPT's column-ordered mask
incidentally aligns with WikiText-2-adjacent tasks when calibrated on the
same distribution; OBS-cancel-block's cross-weight cancellation generalises
better to out-of-distribution calibration data.

Wanda and RIA are insensitive to calibration choice (PPL 2{,}105--4{,}051 in
both conditions), confirming that calibration data only affects methods that
exploit activation statistics via weight correction.

% -----------------------------------------------------------------------
\subsection{Block Size Sweep}
\label{sec:block_sweep}

\begin{table}[htbp]
  \centering
  \caption{%
    Block size sweep on the 1B transformer at 50\% sparsity.
    Dense baseline PPL: 19.5.
  }
  \label{tab:block}
  \small
  \begin{tabular}{lrr}
    \toprule
    Block size $b$ & PPL$\downarrow$ & $\Delta$ from dense \\
    \midrule
    64                                       & 26.38 & +6.86 \\
    \textbf{128 (OBS-cancel-block, ours)}    & \textbf{25.47} & \textbf{+5.95} \\
    256                                      & 25.24 & +5.72 \\
    512                                      & 24.85 & +5.33 \\
    $\infty$ (global OBS-cancel)             & 24.10 & +4.58 \\
    \bottomrule
  \end{tabular}
\end{table}

PPL improves monotonically with block size, confirming that larger blocks
capture more cross-weight cancellation. Gains diminish with each doubling
($128\to256$: $-0.23$; $256\to512$: $-0.39$; $512\to\infty$: $-0.75$),
consistent with most cancellation structure being local within 128-column
windows. Block size 128 is the largest setting that remains numerically
stable across all model sizes; the global variant ($b\to\infty$) is only
feasible on the 1B model.

% -----------------------------------------------------------------------
\subsection{Sparsity Sweep: 1B Transformer}
\label{sec:sweep_1b}

\begin{table}[htbp]
  \centering
  \caption{%
    Sparsity sweep on the 1B transformer. PPL on WikiText-2; dense PPL: 19.5.
    OBS-cancel-block uses the global variant on this model
    (stable at 1B scale; see Section~\ref{sec:block_sweep}).
  }
  \label{tab:sweep_1b}
  \small
  \begin{tabular}{lrrrrrr}
    \toprule
    Sparsity & OBS-cancel-block & SparseGPT & Greedy+corr & Wanda & RIA & AWP \\
    \midrule
    30\% & \textbf{20.2}     & 20.5      & 32.6        & 34.8       & 34.7       & 33.5       \\
    40\% & \textbf{21.2}     & 22.4      & 84.3        & 189.6      & 138.4      & 64.7       \\
    50\% & \textbf{24.1}     & 29.6      & 615         & 3{,}545    & 2{,}105    & 302        \\
    60\% & \textbf{47.0}     & 75.5      & 3{,}507     & 11{,}319   & 10{,}313   & 4{,}371    \\
    70\% & \textbf{4{,}365}  & 7{,}591   & 9{,}065     & 24{,}670   & 11{,}897   & 7{,}447    \\
    80\% & \textbf{14{,}177} & 25{,}929  & 17{,}077    & 11{,}076   & 18{,}369   & 14{,}960   \\
    \bottomrule
  \end{tabular}
\end{table}

\begin{table}[htbp]
  \centering
  \caption{%
    1B transformer at 80\% sparsity: task accuracy.
    All methods collapse to near-random; LAMBADA accuracy is 0 across the board.
  }
  \label{tab:1b_80}
  \small
  \begin{tabular}{lrrrrrrr}
    \toprule
    Method & PPL$\downarrow$ & ARC-e & ARC-c & Hella. & PIQA & Wino. & Avg$\uparrow$ \\
    \midrule
    Dense baseline              & 19.5          & .610 & .268 & .366 & .683 & .525 & .472 \\
    \midrule
    Wanda                       & 11{,}076      & .249 & .219 & .258 & .528 & .494 & .291 \\
    RIA                         & 18{,}369      & .255 & .224 & .258 & .528 & .479 & .291 \\
    SparseGPT                   & 25{,}929      & .253 & .223 & .256 & .537 & .471 & .290 \\
    AWP                         & 14{,}960      & ---  & ---  & ---  & ---  & ---  & ---  \\
    \textbf{OBS-cancel-block}   & \textbf{14{,}177} & \textbf{.266} & \textbf{.224} & \textbf{.259} & .532 & \textbf{.510} & \textbf{.298} \\
    \bottomrule
  \end{tabular}
\end{table}

Table~\ref{tab:sweep_1b} shows PPL across 30--80\% sparsity.
OBS-cancel-block outperforms all methods at every level on the 1B transformer,
with the margin over SparseGPT growing monotonically: $1.02\times$ at 30\%,
$1.23\times$ at 50\%, $1.60\times$ at 60\%, $1.83\times$ at 80\%.
AWP occupies the middle tier: better than Wanda and RIA at $\geq$40\% sparsity
(302 vs.\ 3{,}545 at 50\%), but far behind OBS-cancel-block (24.1).
At 80\% both OBS-cancel-block and AWP are comparable (14{,}177 vs.\ 14{,}960),
yet OBS-cancel-block achieves this in one pass while AWP requires up to
200 iterations per layer---approximately $200\times$ more compute.
Table~\ref{tab:1b_80} adds task accuracy at 80\% sparsity; OBS-cancel-block
retains the best PPL and the best average task accuracy despite all methods
collapsing toward near-random at this extreme compression level.

% -----------------------------------------------------------------------
\subsection{LLaMA-7B}
\label{sec:llama}

\begin{table}[htbp]
  \centering
  \caption{LLaMA-7B at 50\% sparsity. Dense PPL: 8.64. Task accuracies for AWP pending.}
  \label{tab:llama_50}
  \small
  \begin{tabular}{lrrrrrrr}
    \toprule
    Method & PPL$\downarrow$ & ARC-e & ARC-c & Hella. & PIQA & Wino. & Avg$\uparrow$ \\
    \midrule
    Dense baseline            & 8.64  & .723 & .370 & .526 & .749 & .675 & .608 \\
    \midrule
    \textbf{AWP}              & \textbf{10.62} & ---  & ---  & ---  & ---  & ---  & ---  \\
    RIA                       & 11.22 & .601 & .332 & .624 & .732 & .649 & \textbf{.588} \\
    Wanda                     & 11.43 & .602 & .347 & .628 & .732 & .663 & .595 \\
    OBS-cancel-block          & 11.92 & .632 & .323 & .460 & .712 & .643 & .554 \\
    SparseGPT                 & 12.70 & .659 & .349 & .478 & .725 & .643 & .571 \\
    \bottomrule
  \end{tabular}
\end{table}

\begin{table}[htbp]
  \centering
  \caption{LLaMA-7B at 80\% sparsity. Dense PPL: 8.64. Task accuracies for AWP pending.}
  \label{tab:llama_80}
  \small
  \begin{tabular}{lrrrrrrr}
    \toprule
    Method & PPL$\downarrow$ & ARC-e & ARC-c & Hella. & PIQA & Wino. & Avg$\uparrow$ \\
    \midrule
    Dense baseline            & 8.64      & .723 & .370 & .526 & .749 & .675 & .608 \\
    \midrule
    \textbf{AWP}              & \textbf{230.1} & ---  & ---  & ---  & ---  & ---  & ---  \\
    OBS-cancel-block          & 948.5     & .252 & .213 & .260 & .521 & .500 & \textbf{.349} \\
    SparseGPT                 & 1{,}103.6 & .269 & .212 & .258 & .532 & .505 & .355 \\
    Wanda                     & 1{,}252.3 & .253 & .267 & .265 & .499 & .491 & .355 \\
    RIA                       & 1{,}609.5 & .268 & .268 & .261 & .495 & .496 & .357 \\
    \bottomrule
  \end{tabular}
\end{table}

\begin{table}[htbp]
  \centering
  \caption{LLaMA-7B sparsity sweep. Dense PPL: 8.64.}
  \label{tab:llama_sweep}
  \small
  \begin{tabular}{lrrrrr}
    \toprule
    Sparsity & AWP & OBS-cancel-block & SparseGPT & Wanda & RIA \\
    \midrule
    30\% & \textbf{8.92}  & 9.12           & 9.18           & 9.04           & 8.96           \\
    40\% & \textbf{9.44}  & 9.97           & 10.24          & 9.76           & 9.54           \\
    50\% & \textbf{10.62} & 11.92          & 12.70          & 11.43          & 11.20          \\
    60\% & \textbf{14.69} & 18.58          & 20.76          & 18.11          & 18.54          \\
    70\% & \textbf{39.22} & 67.73          & 71.20          & 82.1           & 104.1          \\
    80\% & \textbf{230.1} & 948.5          & 1{,}103.6      & 1{,}252.3      & 1{,}609.5      \\
    \bottomrule
  \end{tabular}
\end{table}

Tables~\ref{tab:llama_50}--\ref{tab:llama_sweep} show results on LLaMA-7B.
Three findings stand out at 7B scale.

\paragraph{AWP is the best method at all sparsity levels on LLaMA-7B.}
AWP achieves the lowest PPL at every sparsity from 30\% to 80\%, with the
margin growing dramatically at high compression: 10.62 vs.\ 11.92
(OBS-cancel-block) at 50\%, and 230 vs.\ 948 ($4.1\times$) at 80\%.
At 7B scale, the optimisation landscape has sufficient redundancy that
iterative IHT---up to 200 gradient-projection steps per layer---finds
substantially better sparse solutions than any one-shot method.
This gain comes at cost: AWP requires approximately $200\times$ more compute
than OBS-cancel-block per layer. For GPU-hosted 7B models where pruning is
a one-time offline operation, this tradeoff may be acceptable. For
neuromorphic deployment at sub-2B scale, Section~\ref{sec:hgrn} shows the
tradeoff reverses entirely: AWP's extra search finds nothing---it collapses
alongside uncorrected methods while OBS-cancel-block succeeds.

\paragraph{Scale determines whether correction is necessary.}
No-correction methods (Wanda, RIA) are competitive at 30--50\% sparsity on
LLaMA-7B, collapsing only beyond 60\% (RIA: $18.5 \to 104 \to 1{,}610$).
Both collapsed to PPL\,$>$\,2{,}000 at 50\% on the 1B model. At 7B scale,
sufficient model redundancy exists that activation-weighted scoring alone
identifies genuinely low-impact weights. OBS-cancel-block and AWP degrade
more gracefully to 948 and 230 respectively.

\paragraph{OBS-cancel-block consistently beats SparseGPT.}
OBS-cancel-block outperforms SparseGPT at every sparsity level on LLaMA-7B
($1.007\times$ at 30\% to $1.163\times$ at 80\%), consistent with the 1B
results. The global OBS-cancel variant fails on LLaMA-7B due to ordering
mismatch and numerical drift as described in Section~\ref{sec:method_ours};
the block-level design resolves both.

% -----------------------------------------------------------------------
\subsection{HGRN-1.3B: Architecture Generalisation}
\label{sec:hgrn}

\begin{table}[htbp]
  \centering
  \caption{HGRN-1.3B at 50\% sparsity. Dense PPL: 14.18.}
  \label{tab:hgrn_50}
  \small
  \begin{tabular}{lrrrrrrrr}
    \toprule
    Method & PPL$\downarrow$ & ARC-e & ARC-c & Hella. & PIQA & Wino. & LMBD & Avg$\uparrow$ \\
    \midrule
    Dense baseline            & 14.18 & .510 & .275 & .480 & .712 & .528 & .383 & .481 \\
    \midrule
    Wanda                     & 404   & .312 & .263 & .295 & .545 & .517 & .003 & .323 \\
    AWP                       & 527   & ---  & ---  & ---  & ---  & ---  & ---  & ---  \\
    RIA                       & 410   & .305 & .255 & .298 & .542 & .511 & .004 & .319 \\
    SparseGPT                 & 20.3  & .467 & .264 & .434 & .676 & .519 & .215 & .429 \\
    \textbf{OBS-cancel-block} & \textbf{19.0} & .461 & .261 & .427 & .669 & .525 & .230 & \textbf{.429} \\
    \bottomrule
  \end{tabular}
\end{table}

\begin{table}[htbp]
  \centering
  \caption{HGRN-1.3B at 80\% sparsity. Dense PPL: 14.18.}
  \label{tab:hgrn_80}
  \small
  \begin{tabular}{lrrrrrrrr}
    \toprule
    Method & PPL$\downarrow$ & ARC-e & ARC-c & Hella. & PIQA & Wino. & LMBD & Avg$\uparrow$ \\
    \midrule
    Dense baseline            & 14.18     & .510 & .275 & .480 & .712 & .528 & .383 & .481 \\
    \midrule
    Wanda                     & 75{,}620  & .249 & .235 & .260 & .513 & .519 & .000 & .296 \\
    AWP                       & 17{,}756  & ---  & ---  & ---  & ---  & ---  & ---  & ---  \\
    RIA                       & 26{,}817  & .255 & .230 & .259 & .516 & .498 & .000 & .293 \\
    SparseGPT                 &  2{,}811  & .269 & .220 & .260 & .533 & .519 & .000 & .300 \\
    \textbf{OBS-cancel-block} & \textbf{1{,}952} & \textbf{.269} & \textbf{.222} & .258 & .527 & .491 & .000 & .294 \\
    \bottomrule
  \end{tabular}
\end{table}

\begin{table}[htbp]
  \centering
  \caption{HGRN-1.3B sparsity sweep. Dense PPL: 14.18.}
  \label{tab:hgrn_sweep}
  \small
  \begin{tabular}{lrrrrr}
    \toprule
    Sparsity & OBS-cancel-block & SparseGPT & AWP & Wanda & RIA \\
    \midrule
    30\% & \textbf{14.92}  & 15.02  & 25.6      & 31.64     & 30.72     \\
    40\% & \textbf{16.06}  & 16.46  & 59.4      & 76.87     & 68.27     \\
    50\% & \textbf{19.02}  & 20.30  & 527       & 584       & 499       \\
    60\% & \textbf{29.14}  & 32.43  & 2{,}761   & 11{,}552  & 11{,}255  \\
    70\% & \textbf{112.7}  & 115.4  & 5{,}123   & 20{,}592  & 14{,}723  \\
    80\% & \textbf{1{,}952}& 2{,}811& 17{,}756  & 75{,}620  & 26{,}817  \\
    \bottomrule
  \end{tabular}
\end{table}

The HGRN results (Tables~\ref{tab:hgrn_50}--\ref{tab:hgrn_sweep}) reproduce
the 1B transformer pattern on a gated recurrent SSM architecture.

\paragraph{Scale, not architecture, determines correction necessity.}
Wanda and RIA collapse at 50\% sparsity (PPL 404, 410) on HGRN-1.3B despite
its comparable parameter count to the 1B transformer. AWP (527) also collapses
at 50\% on HGRN-1.3B, confirming that iterative IHT does not overcome the
fundamental reconstruction error problem at 1B scale: without second-order
weight correction, no-correction and iterative methods alike face the same
barrier. This rules out architecture type as the determining factor: the same
fragility at the 1B--1.3B scale contrasts with the robustness of AWP at 7B.

\paragraph{Architecture-agnostic gains.}
OBS-cancel-block outperforms all other methods at every sparsity level on
HGRN-1.3B ($1.007\times$ at 30\%, $1.068\times$ at 50\%, $1.441\times$ at
80\%). The method operates purely on the weight matrix and its associated
Hessian, with no assumptions about the surrounding computation graph,
explaining its generality across architectures. The contrast with AWP's
behaviour (best at 7B, worst at 1.3B among methods that use covariance
information) underscores that second-order correction is the key ingredient at
smaller scales.


% =====================================================================
\section{Discussion}
\label{sec:discussion}
% =====================================================================

\subsection{Why the Phase Transition Occurs}

The key experimental finding is a scale-dependent inversion: OBS-cancel-block
dominates at 1B--1.3B scale; AWP dominates at 7B scale. We offer the following
mechanistic account.

At 1B--1.3B scale, the model has limited parameter redundancy. Pruning at
50\%+ sparsity removes a large fraction of the effective capacity, and the
uncompensated reconstruction error accumulates across layers until it exceeds
the model's tolerance. In this regime, \emph{the magnitude of the reconstruction
error is the binding constraint}, not the quality of the mask selection.
OBS-cancel-block applies exact second-order weight correction after each
selection step; AWP applies no correction at all (the IHT gradient step adjusts
the surviving weights, but only in the direction of recovering a good initialisation,
not toward the true OBS-corrected minimum). The result is that AWP's mask,
however well-chosen, leaves uncorrected residuals that cascade into perplexity
collapse---the same fate as Wanda and RIA.

At 7B scale, the model has sufficient redundancy that even imperfect weight
correction leaves the network in a recoverable state: no-correction methods like
RIA achieve reasonable PPL at 30--50\%. In this regime, the bottleneck shifts
from \emph{error magnitude} to \emph{mask quality}. AWP's iterative search over
the combinatorial mask space finds sparse configurations substantially better
than any greedy one-shot criterion, including OBS-cancel-block's
cancellation-aware greedy. The benefit of additional search iterations outweighs
the cost of uncorrected residuals.

This account predicts the crossover should occur somewhere between 1.3B and 7B
parameters, and that it should be sparsity-dependent (higher sparsity shifts the
binding constraint back toward error magnitude even at large scale). The data
are consistent: at 60\%+ sparsity on LLaMA-7B, the gap between AWP and
OBS-cancel-block ($14.7$ vs.\ $18.6$ at 60\%) is large but not as extreme
as the 80\% gap ($230$ vs.\ $948$), suggesting error magnitude reasserts as
the constraint at very high compression even on larger models.

\subsection{Implications for Neuromorphic Deployment}

For Intel Loihi and similar event-driven neuromorphic chips, model scale is
bounded by on-chip synaptic memory---current Loihi generations are best suited
to models in the 1B--2B parameter range. This places the target squarely in
the regime where OBS-cancel-block is the optimal choice: better quality than
all baselines, and ${\sim}200\times$ cheaper to compute than AWP.

The architecture-agnostic behaviour of OBS-cancel-block (identical gains on
attention transformers and gated recurrent SSMs) is practically important
because the two dominant architecture families for neuromorphic deployment are
precisely attention-based models (efficient prefill) and SSMs (efficient decode
due to constant-time recurrence). The method works equally well on both without
any architectural adaptation.

\subsection{Limitations and Future Work}

\paragraph{AWP task accuracy.}
We report AWP task accuracy as pending for the LLaMA-7B detailed tables
(Tables~\ref{tab:llama_50} and~\ref{tab:llama_80}). The PPL improvements at
7B scale suggest AWP's task accuracy would also improve, but we do not have
these numbers and do not speculate about their values.

\paragraph{Intermediate scale.}
The crossover between OBS-cancel-block and AWP likely occurs between 1.3B and
7B parameters. Identifying the precise crossover scale and its dependence on
architecture and sparsity is a natural next step, requiring experiments at
2B--4B scale.

\paragraph{Combining correction with iterative search.}
OBS-cancel-block and AWP address different bottlenecks: correction magnitude
and mask quality respectively. A method that applies OBS corrections \emph{within}
an iterative mask search could potentially capture both advantages across all
scales---though the interaction between the correction and the IHT gradient
signal is non-trivial and not explored here.

\paragraph{Fine-tuning after pruning.}
All results in this paper are zero-shot (no fine-tuning after pruning). A small
amount of fine-tuning on the sparse model can recover significant accuracy,
particularly at high sparsity. Whether the relative ordering of methods is
preserved after fine-tuning---especially the OBS-cancel-block vs.\ AWP
inversion across scales---is an open question.

\clearpage
\bibliographystyle{plainnat}
\bibliography{references}

\end{document}
